% Diagram konsep Bab 3: Ekosistem dan batas kerja sama mitra.
% Dibangkitkan tools/render-diagrams.py dari figures/chapter-diagrams.json;
% ubah data JSON lalu jalankan ulang skrip, jangan sunting berkas ini langsung.
\begin{tikzpicture}[
  x=0.92cm, y=0.95cm,
  kotak konsep/.style={draw, line width=0.5pt, rounded corners=2pt,
    align=center, minimum width=2.999cm,
    minimum height=1.330cm, inner sep=0pt, font=\scriptsize,
    text=konsepJudul},
  peran actor/.style={kotak konsep, fill=konsepIsiActor, draw=konsepGarisActor},
  peran service/.style={kotak konsep, fill=konsepIsiService, draw=konsepGarisService},
  peran control/.style={kotak konsep, fill=konsepIsiControl, draw=konsepGarisControl},
  peran data/.style={kotak konsep, fill=konsepIsiData, draw=konsepGarisData},
  peran outcome/.style={kotak konsep, fill=konsepIsiOutcome, draw=konsepGarisOutcome},
  peran risk/.style={kotak konsep, fill=konsepIsiRisk, draw=konsepGarisRisk},
  panah konsep/.style={-{Stealth[length=1.3mm, width=1.2mm]}, line width=0.6pt,
    draw=konsepPanah},
  panah putus/.style={panah konsep, draw=konsepPanahPutus,
    dash pattern=on 2.5pt off 1.5pt},
  label panah/.style={font=\tiny, align=center, fill=white,
    inner sep=1pt, text=konsepLabel},
  catatan/.style={font=\scriptsize, align=left, anchor=north west,
    inner sep=0pt, text width=12.383cm, text=konsepTeks,
    execute at begin node={\hyphenpenalty=10000\hbadness=10000\relax}},
]
  \definecolor{konsepIsiActor}{HTML}{EEF1F6}
  \definecolor{konsepGarisActor}{HTML}{50627D}
  \definecolor{konsepIsiService}{HTML}{EAF2FB}
  \definecolor{konsepGarisService}{HTML}{3F6893}
  \definecolor{konsepIsiControl}{HTML}{E9F3EC}
  \definecolor{konsepGarisControl}{HTML}{316B4B}
  \definecolor{konsepIsiData}{HTML}{F1EDF8}
  \definecolor{konsepGarisData}{HTML}{786394}
  \definecolor{konsepIsiOutcome}{HTML}{E4F2EC}
  \definecolor{konsepGarisOutcome}{HTML}{236C57}
  \definecolor{konsepIsiRisk}{HTML}{FCEDE7}
  \definecolor{konsepGarisRisk}{HTML}{B36947}
  \definecolor{konsepJudul}{HTML}{20394F}
  \definecolor{konsepTeks}{HTML}{455264}
  \definecolor{konsepLabel}{HTML}{35485C}
  \definecolor{konsepPanah}{HTML}{5C7084}
  \definecolor{konsepPanahPutus}{HTML}{A15D40}
  % Panah dan label lebih dahulu; kotak menutupi sisa latar label.
  \draw[panah konsep] (0.93,-2.7) -- (2.77,-2.7);
  \draw[panah konsep] (6.03,-2.7) -- (7.87,-2.7);
  \draw[panah konsep] (4.4,-3.4) -- (4.4,-4.7);
  \draw[panah konsep] (6.03,-5.4) -- (6.2,-5.4) -- (6.2,-3.6) -- (9.5,-3.6) -- (9.5,-3.4);
  \draw[panah konsep] (9.5,-2) -- (9.5,-0.7);
  \draw[panah konsep] (4.4,-0.7) -- (4.4,-2);
  \node[label panah] at (1.85,-2.39) {pembayaran};
  \node[label panah] at (6.95,-2.39) {koordinasi};
  \node[label panah] at (4.4,-4.05) {menetapkan};
  \node[label panah] at (6.2,-4.23) {akses dan batas};
  \node[label panah] at (9.5,-1.35) {layanan pasien};
  \node[label panah] at (4.4,-1.35) {arah kendali};
  \node[peran actor] (klien) at (-0.7,-2.7)
    {\textbf{Perusahaan klien}\\[0.5pt]{\tiny\color{konsepTeks}pembeli paket layanan}};
  \node[peran control] (ork) at (4.4,-2.7)
    {\textbf{Medikara}\\[0.5pt]{\tiny\color{konsepTeks}orkestrator ekosistem}};
  \node[peran actor] (mitra) at (9.5,-2.7)
    {\textbf{Klinik dan}\\\textbf{laboratorium}\\[0.5pt]{\tiny\color{konsepTeks}layanan pelengkap}};
  \node[peran outcome] (pasien) at (9.5,0)
    {\textbf{Pasien}\\[0.5pt]{\tiny\color{konsepTeks}pengguna layanan}};
  \node[peran control] (batas) at (4.4,-5.4)
    {\textbf{Boundary resources}\\[0.5pt]{\tiny\color{konsepTeks}API, sandbox, syarat}};
  \node[peran control] (pilihan) at (4.4,0)
    {\textbf{Pilihan peran}\\[0.5pt]{\tiny\color{konsepTeks}bangun, gabung, mitra}};
  \node[catatan] at (-2.33,-6.47)
    {API dan dokumentasi membuka akses; syarat serta tata kelola membatasi penggunaan kemampuan bersama.};
\end{tikzpicture}
